% !TEX root = final_report.tex
\documentclass[final_report.tex]{subfiles}
\setcounter{chapter}{4}

\begin{document}
\ourchapter{Downstream evaluation}\label{ch:downstream}

Everything in the previous chapter was measured on the pretext task the models were pretrained on. Everything here is measured on new, downstream tasks and the question becomes whether \dphi{} is worth learning (not if it is possible to learn). The models used for finetuning are the models featured in \cref{tab:matched-arms}. There are two tasks. The first is ImageNet-100 classification~\cite{imagenet} - we now add the labels back to the training set and will judge how the model classifies images into those labels. The second task is Horizon Lines in the Wild~\cite{hlw} horizon line estimation. The first task is there to measure the cost of modelling rotation and the second one is there a potential benefit of supervising rotation would appear. Furthermore, the HLW task has a build in specificity control - the task asks for two channels: $\theta$ and $\rho$. $\theta$  is the in-plane orientation whereas $\rho$ is translation-like. A welcome result would be improved performance on $\theta$ while not degrading performance on $\rho$.

Note that finetuning will be performed by full finetune as opposed to a linear-probe as this was the decision made in the upstream~\cite{part}. All statistical protocols were decided upon before seeing any results to protect them from bias.

Based on the held-out HLW split, three results were obtained:
\begin{enumerate}
\item Rotated pixels cost orientation (\texttt{ch2} vs baseline).
\item RoPART$\pm90^\circ$ costs orientation (compared to baseline).
\item RoPART$\pm30^\circ$ beats the baseline but the effect if very small and the confidence intervals graze zero.
\end{enumerate}

We also established in \ref{ch:pretraining} that wider rotation bound was worse than the narrower one on every structural gate. This could have suggested similar pattern to emerge downstream but it appears not to be the case - on classification, $\pm90^\circ$ scores $69.74$ against $\pm30^\circ$'s $70.00$ - there is virtually no difference. So the range cost is confined to the pretext task and is a poor predictor of transfer.


\section{Protocol: end-to-end finetuning}\label{sec:finetune-protocol}
PART~\cite{part} made the choice to finetune the entire model rather than linearly probe. Its downstream results (COCO detection, CIFAR-100, PhysioNet sleep staging) are end-to-end finetunes. Quoting directly:
\begin{quote}
``Once the model is pretrained, we tune the network end-to-end using labeled data in a
supervised setup. Following the standard ViT setup, we eliminate the relative encoder and
substitute it with a linear classification or detection head after the \texttt{[CLS]}
token, which aggregates global input information. Unlike pretraining, we incorporate
randomly initialized learnable position embeddings\ldots{}''~\cite[\S4]{part}
\end{quote}

We followed suit. In our case, preparing the pretrained models involved dropping the \texttt{head} as the relative encoder no longer has a role to play, dropping \texttt{clf} and reinitialising it as the old and new shapes do not match, and re-initialising \texttt{pos\_embed}. Then we finetune according to the recipe in \cref{tab:finetune-recipe} which is identical for all arms and both downstream tasks. Note that nothing is frozen.

\begin{table}[htbp]
  \centering
  \caption[The finetune recipe]{The finetuning recipe}
  \label{tab:finetune-recipe}
  \footnotesize
  \begin{tabular}{@{}l l@{}}
    \toprule
    Setting & Value \\
    \midrule
    backbone            & ViT-B/32 at $224$ px, $87{,}532{,}132$ trainable parameters, none frozen \\
    initialisation      & \py{--init-from <pretrain checkpoint>}; \py{head.*} and \py{clf.*} dropped, \\
                        & \py{pos\_embed} re-initialised (\py{trunc\_normal\_}, std $0.02$) \\
    epochs              & 100, warm-up 5, \py{--eval-every 1} \\
    optimiser           & AdamW, weight decay $0.05$, $\beta = (0.9,\, 0.999)$ \\
    learning rate       & $1 \times 10^{-4}$, cosine to $1 \times 10^{-6}$ \\
    batch size          & 192 \\
    precision           & fp16 - the Turing lab boxes emulate bf16 rather than \\
                        & implementing it, so fp16 is the faster of the two there \\
    dataloader workers  & 6 (the boxes have 8 CPUs) \\
    \bottomrule
  \end{tabular}
\end{table}

One note on reporting: instead of reporting the best epoch (which is a metric biased upwards) we report the mean over the final ten epochs - this applies to both tasks.

\section{Statistical protocol}\label{sec:statistical-protocol}
The statistical protocol was fixed before any training runs finished. This is because the effects measured so far were small and by hand-picking the statistical tests based on whether they give preferred results, we could direct the conclusion in any way we wanted. As a consequence, the held-out test split is scored only once, at the end instead of every epoch. What follows is what was decided:

Throughout the remainder of this section, $e^{(a)}_i$ is the error of arm $a$ on image $i$ of the held-out split, on a channel $e \in \{\theta, \rho, \texttt{horizon}\}$, and $n = 2018$ is the number of images. Each of the three rotation arms is compared to the baseline, never to each other. Note that the HLW held-out errors are from a single scoring run of the 99th epoch weights in contrast to \texttt{val} numbers which use a last-10-epoch mean as stated in \cref{sec:finetune-protocol}.

\begin{description}\setlength{\itemsep}{0.6em}

\item[The statistic: paired per-image difference.] Both arms are scored on the same
images, so scene difficulty cancels instead of entering the noise. $d_i$ is one image's
difference and $\bar d$ the mean; negative means the arm beats the baseline.
\begin{equation}\label{eq:paired-diff}
  d_i \;=\; e^{(a)}_i - e^{(\mathrm{base})}_i,
  \qquad
  \bar d \;=\; \frac{1}{n}\sum_{i=1}^{n} d_i .
\end{equation}

\item[Standard error] $s$ is the spread of the
per-image errors and $s_d$ the spread of their \emph{differences}. The unpaired form carries
the heavy tail of hard scenes ($s \approx 2.13^\circ$ against a mean of $0.97^\circ$) and
gives $0.066^\circ$; the paired form cancels that component and gives $\approx 0.016^\circ$, roughly $4\times$ narrower.
\begin{align}
  \widehat{\mathrm{SE}}_{\text{unpaired}}
    &= \sqrt{\frac{s_a^2 + s_{\mathrm{base}}^2}{n}} \;\approx\; \sqrt{2}\,\frac{s}{\sqrt{n}},
  \label{eq:se-unpaired}\\[2pt]
  \widehat{\mathrm{SE}}_{\text{paired}}
    &= \frac{s_d}{\sqrt{n}},
  \qquad
  s_d^2 = \frac{1}{n-1}\sum_{i=1}^{n} \bigl(d_i - \bar d\bigr)^2 .
  \label{eq:se-paired}
\end{align}
Note that the standard error presented here is only used to justify per-pair approach rather than contribute numerically to any tests.

\item[Paired bootstrap - an interval on the mean difference.] Image indices are resampled
with replacement $B$ times; $\bar d^{*(b)}$ is the mean difference on the $b$-th resample and
$q_p$ its $p$-th percentile. No distributional assumption is made.
\begin{equation}\label{eq:bootstrap-ci}
  \mathrm{CI}_{95}
  \;=\; \Bigl[\, q_{2.5}\bigl(\{\bar d^{*(b)}\}_{b=1}^{B}\bigr),\;\;
                 q_{97.5}\bigl(\{\bar d^{*(b)}\}_{b=1}^{B}\bigr) \Bigr],
  \qquad B = 20{,}000 .
\end{equation}


\item[Wilcoxon signed-rank] Two-sided, on the
same $d_i$, so it is not dominated by the minority of hard scenes that drive $\bar d$. $R_i$
is the rank of $|d_i|$ among the $m$ non-zero differences and $t$ runs over tie group sizes;
$z$ is the tie-corrected normal approximation.
\begin{equation}\label{eq:wilcoxon}
  W^{+} = \!\!\sum_{i \,:\, d_i > 0}\!\! R_i ,
  \qquad
  z = \frac{W^{+} - \tfrac{1}{4}m(m+1)}
           {\sqrt{\tfrac{1}{24}m(m+1)(2m+1) - \tfrac{1}{48}\sum_{t}\bigl(t^3 - t\bigr)}} .
\end{equation}

\item[The corrected threshold: Bonferroni.] Three arms are compared to the baseline on
the primary channel, so $m_{\text{comp}} = 3$.
\begin{equation}\label{eq:bonferroni}
  \alpha' \;=\; \frac{\alpha}{m_{\text{comp}}} \;=\; \frac{0.05}{3} \;=\; 0.0167 .
\end{equation}

\item[Specificity criterion] An effect genuinely about
in-plane orientation must move $\theta$ and leave the translation-like $\rho$ alone. An arm
improving $\rho$ as much as $\theta$ has simply become a better model; one improving only
$\rho$ is evidence \emph{against} the mechanism.
\begin{equation}\label{eq:specificity}
  p_\theta < \alpha
  \quad\wedge\quad
  p_\rho \geq \alpha .
\end{equation}

\end{description}

The bootstrap and the signed-rank test can disagree, and on one comparison they do
(\cref{sec:hlw-results}). Both are reported; neither is suppressed in favour of the
friendlier one.

Note that there are three sources of the noise. The first comes from image sampling as some examples are simply more difficult. To alleviate that, pairing was used as it was cost-free. The second comes from finetune stochasticity. Every arm reported in this chapter is a single finetune seed, so this source is not controlled in the results. A five-seed replication of \texttt{base}, \texttt{ch2} and \texttt{w30} was launched under the protocol above; it had not completed at the time of writing. The third source of noise is pretraining stochasticity. A multiple-seed replication would take 30 hours per arm and was deemed out of budget. The scoring pipeline was first tested on \texttt{val} to rule out wiring issues and to not use \texttt{test}.

\section{Classification on ImageNet-100}\label{sec:in100}
\stub{The no-degradation check on a basic task. Rotated pixels cost accuracy;
supervising \dphi{} pays most of it back, returning to parity with the baseline.}

\section{The orientation-sensitive task}\label{sec:hlw-task}
\stub{Horizon line estimation on HLW. The two prediction channels, and why one is
in-plane orientation while the other is translation-like --- a specificity control
built into the task rather than bolted on.}

\section{Held-out results}\label{sec:hlw-results}
\stub{The two significant deficits, and the RoPART advantage with its specificity
control firing. Reported at the strength the evidence supports: the deficits are
findings, the advantage is a positive signal awaiting a second seed.}

\section{Why the published metric cannot see the effect}\label{sec:auc-blindness}
\stub{The published AUC is dominated by the translation-like channel, so it is blind
to the orientation effect and its arm ordering is unstable between splits. The case
for headlining angular error instead.}

\section{Data integrity}\label{sec:data-integrity}
\stub{What was verified before the numbers were trusted, and the parsing bug that
was caught by a constants check rather than by the results looking wrong.}

\section{Synthesis}\label{sec:synthesis}
\stub{The two downstream measurements against the pretext ordering. Pretext gate
quality turns out not to be a proxy for downstream utility, which qualifies the
range-cost result of \cref{sec:matched-budget}.}

\section{Threats to validity revisited}\label{sec:threats-revisited}
\stub{The interpolation confound and orientation recoverability, revisited against
everything now measured --- what has been closed, and what remains open.}

\ifSubfilesClassLoaded{\pagebreak\printbibliography}{}
\end{document}
